
%======================================================================
\section{v8: mixed channels, hidden pairs, and a Duhamel lemma}
\label{sec:v8mixed}
%======================================================================

\subsection{Feynman densities of channel pairs}

For a bubble whose two lines carry channels \(s\) and \(t\), let
\(x\) be the Schwinger fraction of the \(t\)-line. The identity
\[
 \frac1{l^2(l-p)^2}=\int_0^1\frac{\dd x}{\bigl((l-xp)^2+x(1-x)P^2\bigr)^2}
\]
turns the pole of the slice numerator into a polynomial density
\(\rho_{(s,t)}(x)\). By construction,
\(\int_0^1\rho_{(s,t)}\,\dd x\) equals the circle pole of that pair.
This identity was checked exactly for every bubble pair of every
graph (\loc{slice_D*.json}, field \loc{rho_integrates_to_circle}).
The cone value of the pair is
\[
 \int_0^1\rho_{(s,t)}(x)\,C_{s,t}(x)\,\dd x,\qquad
 C_{s,t}(x)=\ApTr\bigl[e^{-(1-x)A_s}e^{-xA_t}\bigr].
\]
For \(s=t\) this is the circle pole times \(\ApTr e^{-A_s}\).

\begin{definition}[Hidden mixed pair]
A pair with \(s\ne t\) is \emph{hidden} if its circle pole vanishes
but \(\rho_{(s,t)}\not\equiv0\).
\end{definition}

A selection rule based on totals, \(\text{circle}\ne0\), never sees
a hidden pair. The historical lattice records value each pair as
its circle pole times one channel coefficient, so a hidden pair
receives the value zero.

\subsection{The complete list (executed)}

The slice run \loc{logs/v8/slice/hidden_mixed.json} found hidden
pairs in exactly two graphs.
\begin{itemize}
\item \textbf{D01, 13 keys.} Eight of them, \((1,1)\), \((1,6)\),
\((2,2)\), \((2,8)\), \((6,1)\), \((6,6)\), \((8,2)\), \((8,8)\),
were already quarantined; their pairs \((0,1)\) and \((1,0)\) have
\(\rho\propto(2x-1)\).
The other five are \textbf{new}: \((4,4)\), \((4,7)\), \((5,5)\),
\((7,4)\), \((7,7)\). Their pairs \((0,2)\) and \((2,0)\) have
\[
 \rho^{(4,4)}_{(0,2)}=-\frac{\ii(q_0^2-q_3^2)(6x^2-10x+3)}{8\pi},\qquad
 \rho^{(4,4)}_{(2,0)}=-\frac{\ii(q_0^2-q_3^2)(6x^2-2x-1)}{8\pi}.
\]
The \((5,5)\) densities are \(4\) times these, \((7,7)\) equals
\((4,4)\), and \((4,7)\), \((7,4)\) carry the opposite sign.
\item \textbf{D03, 8 keys.} These are the vector-type keys,
already excluded.
\end{itemize}
D02, D04, D05a, D06, D11, D12 and D13 have no hidden pairs.

\begin{proposition}[Status of the v6 branch difference; executed and proved]
\label{prop:v8hidden}
The new D01 keys enter the \(hh\) sum, because the indices 4, 5
and 7 are the transverse components 11, 12 and 22.
Since \(\rho_{(2,0)}(1-x)=\rho_{(0,2)}(x)\) and
\(C_{2,0}(x)=C_{0,2}(1-x)\) by cyclicity, their total cone value is
\[
 2\int_0^1\rho_{(0,2)}(x)\,C_{0,2}(x)\,\dd x
 =-2\int_0^1R(x)\,C_{0,2}'(x)\,\dd x .
\]
Here \(R(x)=\int_0^x\rho_{(0,2)}\), which is proportional to
\(x(x-1)(2x-3)\).
This total is generically nonzero, and on the CR and F branches it
is built from different channel spectra.
Therefore the printed v6 value
\(-\nu\ka^2P^2f_0f'_0/(4\pi)\) is unchanged as the lattice
evaluation of the declared catalogue. As a statement about the
cone graph, however, it is \CONDITIONAL{} on the vanishing, or the
separate evaluation, of this mixed-channel term.
\end{proposition}

\begin{proof}
The symmetry of the densities is read off the executed list. The
cyclicity follows from the cyclicity of the regulated apex trace.
Integration by parts uses \(R(0)=R(1)=0\), which holds because
\(\int_0^1\rho_{(0,2)}=0\).
\end{proof}

\subsection{A per-mode Duhamel identity}

\begin{lemma}[Duhamel identity for Bessel channels]
\label{lem:v8duhamel}
In one angular mode, let
\(A_\mu=-\partial_r^2-r^{-1}\partial_r+\mu^2r^{-2}\) on its
Friedrichs domain, with \(\mu,\nu>0\). With the bulk-subtracted
mode trace,
\[
 (\mu^2-\nu^2)\int_0^1\Tr\bigl[e^{-xA_\nu}\,r^{-2}\,e^{-(1-x)A_\mu}\bigr]\dd x
 =\frac{\mu-\nu}2 ,
\]
and \(C_{\mu,\nu}'(x)=(\mu^2-\nu^2)\Tr[e^{-(1-x)A_\mu}r^{-2}e^{-xA_\nu}]\).
In particular \(C(0)=\Tr e^{-A_\mu}\) and \(C(1)=\Tr e^{-A_\nu}\)
after the same subtraction.
\end{lemma}

\begin{proof}
Differentiate \(x\mapsto e^{-(1-x)A_\mu}e^{-xA_\nu}\) and use
\(A_\mu-A_\nu=(\mu^2-\nu^2)r^{-2}\). Integrate over \(x\in[0,1]\),
and use the per-mode value \(-\mu/2\) (up to a \(\mu\)-independent
constant) of the subtracted trace of \(e^{-A_\mu}\).
\end{proof}

\noindent\emph{Numerical corroboration.} We used Richardson
extrapolation in the radial cutoff. For \(\mu-\nu=1\) the result
was \(0.49993\) against \(0.5\), and for \(\mu-\nu=0.85\) it was
\(0.84990\) against \(0.85\)
(\loc{duhamel_mode_check2.json}). This is corroboration, not an
enclosure.

In the Cartesian frame, all channels share the angular phase set,
and the Bessel orders are \(\abs{k/\alpha+S}\). The CR branch flips
the \(k=-1\) mode of \(S=\pm2\). Lemma~\ref{lem:v8duhamel} reduces
\(C_{0,2}'\) to a mode sum of \(r^{-2}\)-sandwiched traces.
v8 does \emph{not} evaluate that mode sum, and does not do the
\(r^2\)-weighted D04 assembly. Both remain \OPEN.
